\begin{abstract}
Reliable representation learning on biomolecular graphs remains difficult because repeated message passing can wash out local structural cues while simple identity skips do not explicitly model complementary pathways. This paper studies Cross-Residual Graph Neural Networks (CR-GNN), a compact family of graph classification architectures that reuses intermediate states either within a branch, across parallel branches, or through graph-level summaries passed between sequential blocks. The goal is not to claim that every cross-residual design is universally superior, but to determine when cross-path information exchange is useful relative to both plain stacking and stronger residual reuse. We evaluate plain, residual, and cross-residual variants under a unified stratified 5-fold protocol with validation-based early stopping, learning-rate scheduling, gradient clipping, and TensorBoard-backed analysis logging. The final study covers seven datasets, with PROTEINS, DD, and ENZYMES forming the topic-facing biological core and MUTAG, AIDS, Mutagenicity, and MSRC\_9 serving as broader robustness checks. The resulting empirical picture is selective rather than absolute: residual variants remain the strongest default family, winning four of the seven datasets, whereas cross-residual variants win three datasets and outperform the plain baseline on six of seven. These results position CR-GNN as a practical design space for biomolecular graph representation learning, with dataset-dependent cross-residual gains and a clear path toward future integration with richer omics-derived signals.
\end{abstract}
